\section{Mixed variation as canonical rectangle charge}
\label{sec:rectangle-charges}

The remaining canonical cost is concentrated in
\(\Delta_{12}K^0\).  We record its exact geometric meaning.

For \(r,s\ge1\), define the terminal rectangle
\begin{equation}
 R_{r,s}(u,v)=\one_{u\le r}\one_{v\le s}.
 \label{eq:terminal-rectangle}
\end{equation}
It is a rectangle anchored at \((1,1)\) and terminating at \((r,s)\).

\begin{theorem}[Unique terminal-rectangle expansion]
\label{thm:terminal-rectangle-expansion}
Every finitely supported \(F:\N^2\to\mathbb C\) has the exact expansion
\begin{equation}
 \boxed{
 F(u,v)
 =
 \sum_{r,s\ge1}
 \Delta_{12}F(r,s)R_{r,s}(u,v).}
 \label{eq:terminal-rectangle-expansion}
\end{equation}
The coefficients are unique.  Moreover,
\begin{equation}
 \Delta_{12}R_{r,s}(x,y)=\one_{(x,y)=(r,s)}
 \label{eq:rectangle-unit-charge}
\end{equation}
and hence
\begin{equation}
 \cV_\square(F)
 =
 \sum_{r,s\ge1}|\Delta_{12}F(r,s)|
 \label{eq:rectangle-charge-mass}
\end{equation}
is exactly the \(\ell^1\)-mass of the canonical rectangle charges.
\end{theorem}

\begin{proof}
Finite support permits two successive telescoping sums:
\begin{align*}
 \sum_{r\ge u}\sum_{s\ge v}\Delta_{12}F(r,s)
 &=
 \sum_{r\ge u}\sum_{s\ge v}
 \bigl(
 F(r,s)-F(r+1,s)\\
 &\hspace{3cm}
 -F(r,s+1)+F(r+1,s+1)
 \bigr)\\
 &=F(u,v).
\end{align*}
This is \eqref{eq:terminal-rectangle-expansion}.  Direct substitution gives
\eqref{eq:rectangle-unit-charge}.  Applying \(\Delta_{12}\) to any
terminal-rectangle expansion therefore recovers every coefficient, proving
uniqueness and \eqref{eq:rectangle-charge-mass}.
\end{proof}

\begin{corollary}[Weighted rectangle form of the physical coefficient]
\label{cor:weighted-rectangle-form}
For every physical anchor,
\begin{equation}
 \mathfrak X_{\mathfrak a}
 =
 \varepsilon_{\mathfrak a}
 \sum_{r,s\ge1}
 M(r)M(s)\,
 \Delta_{12}K_{\mathfrak a}^{0}(r,s).
 \label{eq:weighted-rectangle-form}
\end{equation}
Thus the canonical Abel problem is exactly a discrepancy estimate for the
literal terminal-rectangle charges against the weight \(M(r)M(s)\).
\end{corollary}

\begin{remark}
Low mixed variation means that the physical datum admits a sparse
canonical superposition of large terminal rectangles.  Sparse support in
the point basis is a different property and can produce maximal rectangle
charge.
\end{remark}

\subsection{Diagonal downsampling}

The next identity is useful both for the general TPC--71 certificate and
for later analytic completions.

\begin{theorem}[Downsampling contraction]
\label{thm:downsampling-contraction}
Let \(F^{[d]}(x,y)=F(dx,dy)\).  Then
\begin{equation}
 \Delta_{12}F^{[d]}(x,y)
 =
 \sum_{r=dx}^{d(x+1)-1}
 \sum_{s=dy}^{d(y+1)-1}
 \Delta_{12}F(r,s).
 \label{eq:downsampling-block-charge}
\end{equation}
Consequently,
\begin{equation}
 \boxed{
 \cV_\square(F^{[d]})
 \le
 \cV_\square(F).}
 \label{eq:downsampling-contraction}
\end{equation}
\end{theorem}

\begin{proof}
Telescoping over the displayed \(d\)-by-\(d\) block gives
\begin{align*}
 &\sum_{r=dx}^{d(x+1)-1}
 \sum_{s=dy}^{d(y+1)-1}\Delta_{12}F(r,s)\\
 &\qquad =
 F(dx,dy)-F(d(x+1),dy)-F(dx,d(y+1))
 +F(d(x+1),d(y+1)),
\end{align*}
which is \(\Delta_{12}F^{[d]}(x,y)\).  The blocks indexed by
\((x,y)\in\N^2\) are disjoint.  Apply the triangle inequality and sum.
\end{proof}

\begin{remark}
For the canonical physical completion,
\cref{thm:primitive-support-collapse} is much stronger:
\(F^{[d]}\) is identically zero for \(d\ge2\).  The contraction theorem is
needed when a nonprimitive analytic completion is deliberately introduced.
\end{remark}

\subsection{Point and rectangle extremes}

For \(p,q\ge1\), write
\[
 e_{p,q}(u,v)=\one_{(u,v)=(p,q)}.
\]
If \(p,q>1\), a point mass \(ce_{p,q}\) has four nonzero rectangle charges
and
\[
 \cV_\square(ce_{p,q})=4|c|.
\]
By contrast, a constant terminal rectangle \(cR_{U,V}\) has one charge and
\[
 \cV_\square(cR_{U,V})=|c|.
\]
These are exact identities, not asymptotic estimates.  They explain why
filling nonprimitive holes may trade pointwise sparsity for rectangular
coherence.
